\Subhead{1.10.2}{Одномодовый и многомодовый режимы генерации}
Моды лазера взаимодействуют через общий запас инверсии. Для амплитуд $a_j$ и интенсивностей $I_j=|a_j|^2$ насыщение можно описать следующей моделью. \paragraph{Несколько мод и конкуренция.}
Для медленных интенсивностей $I_j=|a_j|^2$ распространена модель
\begin{equation}
 \dot I_j=2I_j\left(\mu_j-\sum_k b_{jk}I_k\right).
\end{equation}
Здесь $\mu_j$ — линейный инкремент поля моды, $b_{jj}>0$ — собственное насыщение, $b_{jk}$ при $j\ne k$ — перекрёстное насыщение. Множитель два связан с переходом от амплитуды к интенсивности. В двухмодовом случае совместное решение определяется $b_{11}I_1+b_{12}I_2=\mu_1$, $b_{21}I_1+b_{22}I_2=\mu_2$. Явно $I_1=(\mu_1b_{22}-\mu_2b_{12})/D$, $I_2=(\mu_2b_{11}-\mu_1b_{21})/D$, $D=b_{11}b_{22}-b_{12}b_{21}\ne0$. При $b_{jj}>0$, $I_1,I_2>0$ оно устойчиво по интенсивностям, если $D=b_{11}b_{22}-b_{12}b_{21}>0$: якобиан имеет отрицательный след и определитель $4I_1I_2D$. Для одномодового решения $I_1=\mu_1/b_{11}$ нужно также проверить, что другая мода не растёт: $\mu_2-b_{21}\mu_1/b_{11}<0$.

Однородное уширение усиливает конкуренцию за общий запас инверсии, но само по себе не гарантирует одномодовость: важны геометрия поля, накачка и селекция. Сильное перекрёстное насыщение способствует вытеснению мод. Пространственное выжигание инверсии стоячей волной и неоднородное уширение могут ослаблять конкуренцию и позволять многомодовую генерацию. Несколько генерирующих мод сами по себе не означают синхронизации их фаз и получения коротких импульсов: для этого нужен отдельный механизм связи мод. Подробное устройство конкретных лазеров относится к разделу~4.
